\clearpage
\section*{2026-09-15: DOF records provide a better starting point for parcel history}
\textit{Agent-drafted research entry.}

DOF publishes transaction-level tax-map histories that can recover predecessor
parcel sets directly. Jacob requested a feasibility review before implementing
spatial reconstruction. Public tables identify transaction dates, New/Dropped/
Affected lots, condominium changes, and supporting deed or survey references.
The source should lead the footprint audit, with map comparisons checking land
coverage and resolving partial sites. This is a recommendation; production
characteristics and estimation weights have not been corrected.

The September 15 capture contains 60 header rows for five development blocks,
representing 25 transactions. Retrieving their complete details returns 71
lot-action rows across six blocks. Transaction metadata are identical within
each repeated header group and are joined many-to-one to lot actions. Two
seven-row pages reproduce the corresponding rows of the full ordered query;
all five detail-table counts agree with the service. The citywide header count
is 78,189 rows, with dates from May 20, 2008 to September 14, 2026. Retrieval
needs no login or API key. A full extraction and coverage audit remain to be done.

Wyckoff's May 7, 2025 transaction explicitly retains lot 19, drops 42/51,
and creates 20/21/57. It recovers the three old parcels without a spatial
matching threshold. Third Avenue's December 18, 2025 transaction retains
27/31/35 and drops both 30 and \textbf{135}. The historical tax-map inset
shows lot 135 as a narrow Third Avenue frontage strip; 2023 PLUTO records
108 square feet. The documented five-parcel set sums to \textbf{22,247 square
feet}, versus 19,555 in production. This corrects the September 14 audit's
four-parcel total of 22,139. Only the audit source table and parcel-review
outputs have been updated. Unit counts and parent membership are unchanged.

Jamaica's 2024 condominium creation and 2025 termination identify all six
base lots: 30/65/85/89/94/130. The separate application tracker records a
March 23, 2026 request involving those six plus 98/99 and requesting seven
lots; its status is Returned. These records support the associated parcel
list but do not establish the exact land used by the six residential proposals.
The full associated area of 111,540 square feet must not automatically replace
the current feature. Sullivan's 2021 merger into lot 28 appears in DOF, but
its proposed four-lot split does not appear in the retrieved records.
Bruckner's retrieved current map likewise retains the older 1/34/38 layout.

An Affected lot participated in a tax-map edit; that label alone does not
identify a development's land. Transactions can cover several properties,
and condominium sites can include retained uses. The action tables contain
no historical polygons or complete allocation of each old parcel to each
new one. Proposed subdivisions may also be absent from completed changes.
DOF dates, tracker completion dates, PLUTO APPDates, filing dates, and
developer decision dates are distinct. The recommended next step is to trace
DOF histories for all estimation parents, check reconstructed site coverage
against dated maps, and review the remaining ambiguous footprints.

\small
\textit{Sources and reproduction.} The acquisition audit retains 50 source
responses with exact URLs and SHA-256 hashes in
\path{tasks/audits/download_parent_review_documents/code/dof_history_2026-09-15/}.
The \href{https://services6.arcgis.com/yG5s3afENB5iO9fj/ArcGIS/rest/services/DTM_ETL_DAILY_view/FeatureServer}{public DOF service}
and \href{https://nycdof.maps.arcgis.com/apps/dashboards/4c90c216b3064767805833ff147a7e2a}{application tracker}
are the primary sources. The audit's \texttt{review\_dof\_history.R} produces
\texttt{dof\_lot\_changes.csv} and \texttt{dof\_applications.csv} with SaveData
reports; the detailed assessment is
\path{tasks/audits/audit_hdb_mappluto_condo_recovery/report/dof_history_feasibility.md}.
Existing pilot results remain provisional. No model estimation was run.
\normalsize
